\documentclass[11pt,a4paper]{article}
\usepackage[T1]{fontenc}
\usepackage{lmodern,amsmath,amssymb,amsthm,mathtools,booktabs,microtype,xcolor}
\usepackage[margin=26mm]{geometry}
\usepackage[colorlinks=true,linkcolor=blue!45!black,citecolor=blue!45!black,urlcolor=blue!45!black]{hyperref}
\newtheorem{theorem}{Theorem}[section]
\newtheorem{lemma}[theorem]{Lemma}
\newtheorem{proposition}[theorem]{Proposition}
\theoremstyle{remark}\newtheorem{remark}[theorem]{Remark}
\newcommand{\Q}{\mathbb Q}\newcommand{\Z}{\mathbb Z}\newcommand{\HH}{\mathcal H}
\DeclareMathOperator{\Tr}{Tr}\DeclareMathOperator{\Nm}{Nm}
\setlength{\emergencystretch}{2em}
\title{An explicit quartic in 322 variables}
\author{}\date{9 November 2026}
\begin{document}\maketitle
\begin{abstract}
We specify an integer quartic in 322 total variables whose rational input
projection is the integers under the stated arithmetic premises of the
elliptic construction. The local construction uses four factors for
sixteen tests. Reciprocal abscissas have even valuations at every good
unramified place, allowing smaller powers in those factors. The auxiliary
point needs formal depth ten, which still gives comparison precision
forty. The complete quadratic system, expanded quartic, and independent
exact verification certificates accompany the paper. The displayed
member has 4992 monomials and coefficient height $80989^{109858}+1$.
\end{abstract}

\section{Families and arithmetic premises}
All variable counts include the distinguished input $X$. For a fixed
integer $N\ge1$ set $A_N=80989^{2N+1}$. The coefficient tables define 294 primitive integer quadratic
polynomials and their sum of squares
\[
                     P_N=\sum_{j=1}^{294}Q_{N,j}^2
\]
in 322 coordinates. The parameters $N,A_N$ are fixed coefficients,
not variables.

\begin{theorem}\label{thm:families}
Assume the arithmetic premises specified below. For every $N\ge27464$
and every $x\in\Q$,
\[
 x\in\Z\quad\Longleftrightarrow\quad
 \exists\mathbf y\in\Q^{321}\ P_N(x,\mathbf y)=0.
\]
Thus $A_{27464}=80989^{54929}$ specifies a defining quartic.
Every $P_N$ with $N\ge1$ has degree exactly four, content one, and
4992 nonzero coordinate monomials. For every $N\ge23$,
\[
                         H(P_N)=80989^{4N+2}+1.
\]
The explicit member has coefficient height $80989^{109858}+1$,
whose binary length is 1791283.
\end{theorem}

The following arithmetic inputs have precisely the scope stated in
\cite{Dependencies,Shared,Explicit}. Set
\[
 K=\Q(\rho),\quad \rho^2=2,\quad D=11-6\rho=(3-\rho)^2,
 \qquad E:y^2=x^3-Dx,
\]
\[
 P=2(1+\rho,2\rho)
  =\left(\frac{57-28\rho}{8},\frac{-660+417\rho}{32}\right).
\]
We retain the group assertion $2E(K)=\Z P$ and
$h(x([n]P))\ge n^2-6$ for nonzero integer $n$; the five-contact
height bound with exponent nineteen and its stated exceptional primes;
the pole obstruction outside $2,3,17$ and its simultaneous realization
for inputs with square reduced denominator, the specified local
conditions at those three primes, and real value in $(0,17)$; and the
simultaneous prime-pattern and conjugate reduction-order assertions,
including finite congruence preparation and real approximation.
The quantitative selection and fixed denominator dilation used in the
construction follow from the analytic estimates as detailed in
\cite{Selection}. These finite coefficient calculations do not prove
those arithmetic inputs.

We use the effective sum-height theorem of \cite{Sums}, a proved
height consequence with explicit constants. No contact-height amplification or new
prime-distribution premise is used in this paper.

The retained equations outside the local factors, including the scalar
normalizer, are exactly the equations specified in
\cite{Shared,Explicit}. Their source hashes and a comparison
of every retained row are included in the independent certificates.
The following sections prove every changed local implication, including
its domains and its effect on completeness.

\section{Even valuations and four local factors}
Let $v$ be a place of $K$ above an unramified prime $p>76176$,
normalized by $v(p)=1$. In particular $v(D)=0$. Write $f(Q)=1/x(Q)$
when defined, and put $f(O)=0$, with $v(0)=+\infty$.

\begin{lemma}\label{lem:even}
If $Q\in E(K_v)$ and $x(Q)\ne0$, then every finite value of
$v(f(Q))$ is even.
\end{lemma}
\begin{proof}
Put $a=v(x(Q))$. If $a<0$, then $v(x^3)<v(Dx)$, so
$2v(y)=3a$ and $a$ is even. If $a>0$, then
$v(Dx)<v(x^3)$, so $2v(y)=a$ and again $a$ is even.
For $a=0$ the assertion is immediate. If $y=0$, the nonzero abscissas
are $\pm(3-\rho)$, both units, so this case is included.
\end{proof}

The root points have nonzero integer indices
\[
 Q_1=[n]P,\quad Q_2=[a]P,\quad Q_3=[t]P,\quad Q_4=[b]P.
\]
The auxiliary points and their exact group differences are
\[
 U_0=H,\quad V_0=H-P,\quad U_i=[u_i]P,\quad
 V_i=U_i-Q_i\quad(1\le i\le4).
\]
For $i>0$ the charts allow either $U_i$ or $V_i$ to equal $O$.
The denominator points $H,V_0$ are excluded from the exceptional
classes in Section~\ref{sec:formal} before their ratios are used.
Let $\sigma$ be conjugation and define the tests
\[
 A_i=f(U_i),\qquad B_i=\sigma f(V_i),\qquad D_H=f(H),
\]
\[
 C_i=L_i-L_i',\qquad I=L_1-\beta,\qquad \beta=51X^2+2504,
\]
\[
 G=L_3^2-L_2(L_2-2)L_3L_4+(L_2-2)^2L_4^2-1.
\]
The intended orders are one for the nine $A_i,B_i$, twenty for
$D_H$, ten for $C_1,I$, and forty for $C_2,C_3,C_4,G$.

For each nonzero square scale $a_*$ the detector retains
\[
 t_*=a_*\lambda,\quad \delta=t_*-1,\quad
 F(t_*)=\prod_{\theta\in\Theta}(t_*-\theta),\qquad
 \Theta=\{0,-1/2,-1,1/2399,-32/407\},
\]
\begin{equation}\label{eq:clearedweights}
 F(t_*)e=(a_*\delta)^5,\qquad \delta^3q=e.
\end{equation}
On $\delta\ne0$ these give
\[
 R=\frac{F(t_*)}{a_*^5\delta^6},\qquad q=\frac1{R\delta^4}.
\]
At a contact prime $p$, $\delta$ is a unit and $\nu=v_p(R)$ is odd,
so $v(q)=-\nu$ at every place above $p$.

Use the four factors
\begin{align}
 T_1&=1+q^{79}C_2^2+q^{38}D_H^2+q^4A_1^4+q^8B_0^8,
       \label{eq:factor1}\\
 T_2&=1+q^{79}C_3^2+q^{38}C_1^4+q^4A_2^4+q^8B_2^8,
       \label{eq:factor2}\\
 T_3&=1+q^{79}C_4^2+q^{38}I^4+q^4A_3^4+q^8B_3^8,
       \label{eq:factor3}\\
 T_4&=1+q^{79}G^2+q^2A_4^2+q^4B_1^4+q^8B_4^8,
       \label{eq:factor4}
\end{align}
and let $W=T_1T_2T_3T_4$.

\begin{lemma}\label{lem:classes}
At a contact place, the four finite nonconstant valuations in each
$T_j$ have exact powers of two divisors $1,2,4,8$, respectively.
Consequently $T_j\ne0$, $v(T_j)\le0$, and $v(T_j)=0$ implies
$T_j\equiv1$ in the residue field.
\end{lemma}
\begin{proof}
In the first column of \eqref{eq:factor1}--\eqref{eq:factor4} the
valuation is $2v(C)-79\nu$, which is odd. In the second column it is
one of
\[
 2\bigl(v(D_H)-19\nu\bigr),\qquad
 2\bigl(2v(C)-19\nu\bigr),\qquad
 2\bigl(v(A_4)-\nu\bigr).
\]
The parenthesis is odd, by Lemma~\ref{lem:even} when needed.
The third and fourth column valuations are
$4(v(f)-\nu)$ and $8(v(f)-\nu)$, with even $v(f)$.
They therefore have the asserted exact divisors. The constant term
has valuation zero, distinct from every nonconstant valuation.
Thus the minimum is attained exactly once, proving all assertions.
A zero test removes its term and does not affect the argument.
\end{proof}

\begin{proposition}\label{prop:extraction}
At a contact prime, suppose the retained cleared trace equations and
pole system hold. Then $\nu>0$, and there is one place above that
prime at which all sixteen tests have their intended orders.
\end{proposition}
\begin{proof}
Write $W=w_0+w_1\rho$. The literal trace equations are
\begin{equation}\label{eq:trace}
 w_0^2-2w_0T_0-2w_1^2=0,\qquad
 e\delta\mathfrak b=T_0^2,\qquad e\delta=R^{-1}.
\end{equation}
Lemma~\ref{lem:classes} shows that $W\ne0$. If $w_0=0$, the first
row would imply $w_1=0$, a contradiction. Hence
\[
 T_0=\frac{\Nm(W)}{2w_0}=\frac1{\Tr(W^{-1})}\ne0,
 \qquad \mathfrak b=\frac{R}{\Tr(W^{-1})^2}\ne0.
\]
Put $\ell=\max_{v\mid p}v(W)\le0$. The trace inequality gives
$v_p(\Tr(W^{-1}))\ge-\ell$ and therefore
\[
                v_p(\mathfrak b)\le\nu+2\ell.
\]
This valuation is odd. The pole obstruction forces
$0<v_p(\mathfrak b)\le\nu+2\ell$, whence $\nu>0$ and
$\ell>-\nu/2$. At a place attaining $\ell$, each factor has
valuation at least $\ell$, because all four factor valuations are
nonpositive and sum to $\ell$. Each nonzero term in that factor
therefore has valuation strictly greater than $-\nu/2$.

For the first column this gives
$v(C)>(79/2-1/4)\nu\ge39.25$, hence $v(C)\ge40$.
The second-column $C_1,I$ terms give
$v(C)>(38/4-1/8)\nu\ge9.375$, hence order ten.
For $D_H$ the inequality is
\[
             v(D_H)>(19-1/4)\nu\ge18.75.
\]
Its valuation is even, so it is at least twenty.
The $A_4$ inequality gives $v(A_4)>3\nu/4\ge3/4$;
those for the third and fourth columns give respectively
$v(f)>7\nu/8$ and $v(f)>15\nu/16$.
Their even integer valuations are therefore at least two, proving
all order-one requirements. Zero tests satisfy every required order.
\end{proof}

\paragraph{Degenerate detector parameters.}
The cleared identities \eqref{eq:clearedweights} have exactly the domain
properties proved in \cite[Lemma 2.1]{Shared}. Since $a_*\ne0$ and
$1\notin\Theta$, $F(t_*)=0$ is impossible. If $q=0$, then $e=0$
and $\delta=0$. At $t_*=1$ the contact radical is one, because all
prime divisors of the numbers $1-\theta$ lie below the fixed cutoff.
Thus no contact argument uses a reciprocal on its zero stratum.

\section{Formal comparisons at depth ten}\label{sec:formal}
This section proves the depth implication without presupposing
integrality of the ratios. It also deals with the enlarged $H$ domain
and identity numerators.

The retained chart for $H$ is
\[
 xf=1,\qquad S^2=x-Df,\qquad H=(x,xS).
\]
The group calculation in the elliptic source writes
$H=[m]R_0+T$, where $P=2R_0$, $m\in\Z$, and $T\in E[2](K)$.
Put $h=m/2$ for its rational free index; this notation does not
assert a half-integer multiplication on the curve.

\subsection{Denominator exclusions}
Every point of $E[2](K)$ or $\pm P+E[2](K)$ which is permitted by
the finite chart has reciprocal abscissa a unit at all primes above
$p>76176$. This follows from the explicit translation identities
and norm values
\[
 \Nm(x(P))=41^2/64,\quad
 \Nm(x(P)-(3-\rho))=17^2/64,\quad
 \Nm(x(P)+(3-\rho))=63^2/64
\]
given in the source. If $H$ were such a point,
Proposition~\ref{prop:extraction} would contradict any contact prime.
All comparison radicals would therefore be one. The retained height
inequality for the root labels, together with the scalar lower bound
at the stated coefficient threshold, excludes this possibility exactly
as in \cite{Explicit}. This argument uses no quotient of
logarithms. Consequently $m\notin\{0,2\}$, and $H$ and $V_0=H-P$
are nontorsion.

The denominator logarithm has the same global nonvanishing
certificate as in the source. In particular its rational approximant
has numerator with no root modulo the split prime 71 and denominator
with no root modulo the split prime 7. Their leading coefficients are
units at those primes. A root in $K$ would give an integral local root
and a forbidden residue root. Thus the rational approximant never vanishes on $K$. The retained
chart gives $z(V_0)\ne0$, since $V_0$ is nontorsion and neither its
abscissa nor its ordinate is zero. Their product, the algebraic
approximate logarithm used as the denominator, is therefore nonzero.
The analytic formal logarithm is treated separately below by the
formal logarithm isometry. All unchanged chart denominators are
justified by the retained source equations.

\subsection{Formal logarithm errors}
Write $z=-x/y$. At a nonidentity formal point of depth
$r=v(z)\ge1$, one has $v(f)=2r$ and
\[
 \log_E(Q)=z(Q)\mathcal A(Q),\qquad
 \mathcal A(Q)=G(Df(Q)^2),\qquad
 G(w)=\sum_{j\ge0}g_jw^j,
\]
where
\[
 g_0=1,\quad g_1=-2/5,\quad
 g_j=\frac{4j-5}{4j+1}g_{j-1}\quad(j\ge2).
\]
This is the formal logarithm expansion established in the elliptic
source; it also follows by integrating the invariant differential.
Set $\mathcal A(O)=1$.

\begin{lemma}\label{lem:logerror}
At every place in question,
$\mathcal A(Q)\in1+p^{4r}\mathcal O$ for a nonidentity formal point
of depth $r$. The relative errors of the retained quadratic Taylor
polynomial $G_2$ and rational approximant $\mathcal R$ are at least
$12r$ and $40r$, respectively. At $O$ the errors are zero.
\end{lemma}
\begin{proof}
The recursion gives $v(g_j)\ge-(4j+1)/(p-1)$, since each denominator
divides a product of factors of $(4j+1)!$. When $4j+1<p$, it is
integral. Since $v(Df^2)=4r$, the nonconstant terms have valuation
at least $4r$: this is immediate in the integral range, and for the
remaining $j\ge19044$ follows from
$4rj-(4j+1)/(p-1)\ge4r$.
The same argument starting at $j=3$ or $j=10$ gives the Taylor tail
bounds $12r$ or $40r$.
For the retained rational approximant, the exact identity
\[
 Q(w)G_9(w)-N(w)
 =\frac{2970726}{254443709}w^{10}
 -\frac{41697666}{6361092725}w^{11}
\]
is part of its coefficient certificate. All denominators in the
approximant are supported at primes at most 37, and $Q(0)=1$;
therefore $\mathcal R-G_9$ has valuation at least $40r$.
Division by the unit $\mathcal A(Q)$ preserves these bounds.
At $O$ the analytic unit quotients are defined by their constant
value one, so no expression $\log_E(O)/z(O)$ is used.
\end{proof}

\subsection{Integrality before precision}
Fix the place supplied by Proposition~\ref{prop:extraction}. The
positive $f$ valuations put $H,U_i$ on the primary formal branch,
and $V_0,V_i$ on the conjugate formal branch. In particular
$v(z(H))\ge10$. The group identities
\[
 [m]U_i=[2u_i]H,\qquad
 [m-2]V_i=[2(u_i-n_i)]V_0
\]
kill the two-torsion component of $H$. Applying the formal logarithm
at the respective places gives the full ratios
\[
               r_i=u_i/h,\qquad r_i'=(u_i-n_i)/(h-1).
\]
The denominators are nonzero by the preceding exclusions, and the
formal logarithms of $H,V_0$ are nonzero by the formal logarithm
isometry at these places \cite[Theorem IV.6.4]{Silverman}.

The primary approximate ratio is exactly
\[
 L_i=r_i\frac{\mathcal A(H)}{\mathcal A(U_i)}.
\]
The alternate ratio is $r_i'$ times a conjugated quotient of units
$G_2/\mathcal A$ or $\mathcal R/\mathcal A$, with the denominator
using $\mathcal R$ at $V_0$. This remains true when a numerator point
is $O$, since both the approximate ratio and its rational index ratio
then equal zero.

At the odd prime $p$, $h,h-1\in\Z_{(p)}$ differ by one, so at least
one is a unit. The corresponding full ratio is integral because the
integer labels are integral. Multiplication by the indicated analytic
units preserves its valuation. The positive-order test $C_i=L_i-L_i'$
then transfers integrality to the other approximate ratio and hence
to both full ratios. This step uses only the initially known formal
branches and the unit assertion in Lemma~\ref{lem:logerror}; it does
not use the deeper precision conclusion.

Now $L_i=z(U_i)/z(H)$ is integral. If $U_i\ne O$, it follows that
$v(z(U_i))\ge v(z(H))\ge10$. The primary relative error is therefore
at least forty. If $U_i=O$, its error is exactly zero. The alternate
relative error is at least twelve for $i=1$ and forty for $i>1$.
The full ratios are integral, so these relative errors give absolute
congruences. Together with the cross tests, they yield
\[
 L_i\equiv r_i\equiv r_i'\pmod{p^{k_i}},
 \qquad k_1=10,\quad k_i=40\ (i>1).
\]
Multiplying by the integral numbers $h$ and $h-1$ and subtracting
shows
\begin{equation}\label{eq:indices}
           L_1\equiv n\pmod{p^{10}},\qquad
           L_i\equiv n_i\pmod{p^{40}}\quad(i>1).
\end{equation}
At an inert prime conjugation preserves the local valuation, and
the identical argument applies. This proves the required local
comparison without an assumption that either numerator is nonzero
or that $h$ is an integer.

\section{Global soundness and completeness}
Equation~\eqref{eq:indices}, the order-ten input test, and the
order-forty Lucas test give precisely the divisibility statements
used by the retained global height proof. The Lucas polynomial has
integer coefficients and all three relevant approximate ratios are
integral at the extracted place, so substitution of the integer
indices changes its value by a multiple of $p^{40}$. The input test
similarly gives the original order-ten divisibility for the numerator
of $n-\beta$. Thus the same height contradictions force the Lucas
identity and then $n=\beta$. Since $51$ is squarefree,
$51x^2+2504\in\Z$ implies $x\in\Z$: a reduced denominator $v$ would
satisfy $v^2\mid51$. The scalar lower bound and every numerical
threshold in the explicit construction are unchanged. This proves soundness at the scope stated in
Theorem~\ref{thm:families}.

For completeness retain the source's choice of integer root labels,
scalar point, scales, fixed local preparation, contact primes, and
common auxiliary group index. At each selected primary place its
construction gives depth at least twenty for $H$, formal depth at
least one for each other relevant point, cross precision at least
forty for $C_2,C_3,C_4$, input cross precision at least twelve,
and exact target and Lucas relations to the needed precision.
Consequently every nonconstant term in
\eqref{eq:factor1}--\eqref{eq:factor4} has positive valuation when
$\nu=1$. For example the lower bounds for the first factor are
$1,42,4,8$. Hence $W\equiv1$ at the designated place. At the opposite
place each factor has valuation at most zero and is one modulo the
prime if its valuation is zero. Thus the trace of $W^{-1}$ reduces
to 1 or 2 and is a unit, as required for the retained pole input.

At the fixed exceptional primes, take neighborhoods of the already
prepared points on which all test valuations are bounded below.
For $\lambda=u/(D_0^2w^2)$, choose the fixed integer $D_0$ sufficiently
divisible before choosing the primes. Since the zero contact makes
$u$ a selected prime outside the fixed exceptional set, the argument
of \cite[Section 3]{Shared} gives
\[
        v(q)=2v(a_*)-3v(\lambda)
                \ge2v(a_*)+6v(D_0).
\]
All exponents of $q$ in the new factors are positive, so the new
terms are arbitrarily small simultaneously on those neighborhoods.
This preserves the local squareclasses and lower valuation bounds
at $2,3,17$ and the scalar conditions.

At both real embeddings, $q$ tends to zero through positive values
and the test values remain bounded on the selected real neighborhoods.
Every test power in the four factors is even. Thus $W\to1$ through
positive values and the pole input has the same positive decay as in
the retained construction. Outside the selected primes the detector
has even valuation; the squared trace denominator therefore preserves
the square reduced-denominator condition. The pole realization gives
all remaining witnesses. This proves completeness without strengthening
the simultaneous prime-pattern input.

\section{Quadratic specification and exact verification}
The scalar weight circuit is
\[
 q_2=q^2,\quad q_4=q_2^2,\quad c q_4=q,\quad q_8=q_4^2,
 \quad c q_{19}=q_8^2,\quad q_{38}=q_{19}^2,
 \quad c q_{79}=q_{38}^2.
\]
At a contact these force the displayed powers, and their degenerate
stratum is already covered by \eqref{eq:clearedweights}.
The required squared tests are retained coordinates or affine
combinations of them. The fourth powers of $B_1$ and the eighth
powers of $B_0,B_2,B_3,B_4$ are already present in the logarithm
circuits. Exactly five additional field values are needed:
\[
                  C_1^4,\ I^4,\ A_1^4,\ A_2^4,\ A_3^4.
\]
Each is computed by two rational quadratic equations from its existing
square, using $(u+v\rho)^2=u^2+2v^2+2uv\rho$.

For each detector the four factors cost eight rational coordinates.
Two intermediate field products cost four more; the final field
product uses the retained coordinates of $W$. Thus each detector
costs twelve new coordinates and fourteen quadratic equations.
The complete count is
\[
\begin{array}{lr}\toprule
\text{retained coordinates, including input}&276\\
\text{five shared field powers}&10\\
\text{four factors and intermediate products for three detectors}&36\\\midrule
\text{total coordinates}&322\\\midrule
\text{retained quadratic rows}&242\\
\text{shared power rows}&10\\
\text{factor and product rows}&42\\\midrule
\text{total quadratic rows}&294\\\bottomrule
\end{array}
\]

The exact scripts use rational arithmetic, clear each row's
coefficient denominators, and divide by its integer content. Since a
sum of rational squares vanishes precisely when each term vanishes,
the resulting quartic has the required input projection.
The independent checker imports no code from the producer. It
reconstructs every new row, checks every retained row against the
pinned source, derives the free fourth and eighth powers by arithmetic
in $\Q(\sqrt2)$, and expands the entire sum of squares by unordered
pairs. It also verifies coordinate occurrence, degree, content,
coefficient support, the height comparison, and 70912 supplementary
finite valuation cases. The finite cases supplement, and do not replace,
the valuation proofs above.

The unique coefficient involving $A_N^2$ is the coefficient of the
actual scalar contact product squared, equal to $A_N^2+1$.
Every other coefficient is a nonzero constant or a nonzero constant
multiple of $A_N$. This proves support independence for positive
$A_N$. The maximum constant coefficient $C$ has 1520 binary digits,
and the maximum linear multiplier is
$L=2\cdot80989^2=13118436242$. Exact integer comparison gives
\[
 C>L^2,\qquad 80989^{90}+1<C<80989^{94}+1,\qquad 80989^{47}>L.
\]
Therefore $H(P_N)=\max(C,A_N^2+1)$ for every positive integer $N$,
and $H(P_N)=A_N^2+1$ exactly when $N\ge23$.
This algebraic height threshold does not alter the defining threshold
$N\ge27464$ in the explicit theorem.

\begin{thebibliography}{99}
\bibitem{Dependencies} Project repository,
\emph{Arithmetic dependencies of the defining-polynomial constructions},
\nolinkurl{arithmetic_dependencies_20261108_v8.tex}.
\bibitem{Shared} Project repository,
\emph{Shared local weights and defining quartics},
\nolinkurl{shared_local_weights_20261109_v12.tex}.
\bibitem{Explicit} Project repository,
\emph{An explicit quartic in 342 variables},
\nolinkurl{quartic342_20261109_v12.tex}.
\bibitem{Sums} Project repository,
\emph{Effective heights of sums of elliptic abscissas},
\nolinkurl{effective_elliptic_sums_20261108_v10.tex}.
\bibitem{Selection} Project repository,
\emph{Prime selection with fixed congruences and separated reduction orders},
\nolinkurl{prime_selection_20261109_v12.tex}.
\bibitem{Silverman} J. H. Silverman,
\emph{The Arithmetic of Elliptic Curves}, second edition,
Graduate Texts in Mathematics 106, Springer, 2009.
\end{thebibliography}
\end{document}
